% SPDX-FileCopyrightText: 2026 Yiyu Zhou <yiyu@yiyuzhou.io>
%
% SPDX-License-Identifier: GPL-3.0-or-later

\chapter{Making a cue}
\label{ch:cues}

This chapter creates a show of your own next to the storm show, with a cue
that plays one of its sounds.  Besides audio cues, cinq provides video and
picture cues, fade and stop cues, which act on other cues, and wait and memo
cues; later chapters cover each of them.

\section{A new show}

Close the main window to quit cinq, then start it again.  It opens a new show
with an empty cue list, named \texttt{untitled} until you save it.

\section{Adding a cue}

\begin{enumerate}
\item Press \ui{Add cue} in the toolbar.  The list beside it sets the kind
  of cue to add, \ui{Audio} by default.  The new cue has no file yet, so its
  state reads \texttt{NO FILE}, and the sidebar shows a warning, as in
  Figure~\ref{fig:added}.
\item Press \ui{Choose} next to the \ui{File} field of the inspector, and
  choose \texttt{rain.flac} in the \texttt{sounds} directory of your copy of
  the demo.  The cue takes the file's name, \texttt{rain}, and the warning
  disappears.
\item Press \ui{GO}\@.  The rain plays once, for eight seconds.
\end{enumerate}

\begin{figure}[H]
  \centering
  \includegraphics[width=\textwidth]{screenshots/added.png}
  \caption{New show, with an audio cue that has no file yet}
  \label{fig:added}
\end{figure}

GO skips a cue that cannot play, and clicking its warning selects it.  The
\ui{File} field also accepts a typed path, either absolute or relative to the
show file's directory.

\section{Setting the level}

\begin{enumerate}
\item Click the \ui{Levels} tab of the inspector.
\item Drag the slider to $-20$~dB, or type \texttt{-20} in the field beside
  it.
\item Press \ui{Standby selected}, then \ui{GO}\@.  The rain plays more
  quietly.  Move the slider while it plays: level changes apply immediately.
\end{enumerate}

Levels range from $-60$~dB, which is silent, to $+12$~dB\@.  At the default
of 0~dB, the cue plays its file unchanged.

\section{Looping and trimming}

\begin{enumerate}
\item Click the \ui{Time \& Loops} tab, and check \ui{Loop until stopped}.
\item Press \ui{Standby selected}, then \ui{GO}\@.  The rain restarts each
  time it reaches its end.
\item Press the rain's \ui{Stop} button in the sidebar.
\end{enumerate}

\ui{Start} and \ui{End}, on the same tab, limit playback to part of the file,
in seconds from its beginning.  In the storm show, cue~7 plays the thunder
from 0.2 to 3.5 seconds, and \ui{Duration} shows the resulting length, as in
Figure~\ref{fig:time}.

\begin{figure}[H]
  \centering
  \includegraphics[width=\textwidth]{screenshots/time.png}
  \caption{\ui{Time \& Loops} tab of an audio cue that plays part of its
    file}
  \label{fig:time}
\end{figure}

\section{Editing the list}

The \ui{Basics} tab holds the cue's \ui{Number}, which can be any text, such
as \texttt{1}, \texttt{2.5}, or \texttt{E1}; its \ui{Name}; and its
\ui{Notes}, which the transport shows while the cue is on standby.
\key{Alt+ArrowUp} and \key{Alt+ArrowDown} move the selected cue up and down
the list, and \ui{Delete cue} deletes it.

The status bar at the bottom of the window shows each change along with the
key that undoes it, and after an undo, the key that redoes it.  cinq never
asks for confirmation, since undo is the way back.

Press \key{Control+s}, and save the show as \texttt{my-show.cinq.toml} in
your copy of the demo, where Chapter~\ref{ch:fades} opens it again.

\section*{Exercise}

You can now make a cue that plays a sound once or in a loop, at the level of
your choice.  Add a second audio cue for \texttt{sounds/thunder.wav} that
plays only its first two seconds.
